生成式推荐的效果高度依赖 item tokenizer，因为它决定了每个 item 如何被映射为离散语义 ID。近期已有不少工作意识到 collaborative information 应当进入 tokenizer，但它们大多把协同信号看作一种全局统一表示，再在整个 tokenizer 中统一融合。我们对 MiniOneRec 的系统分析表明，这种 framing 对层级语义 ID 仍然过于粗糙。我们发现，在最终 \texttt{sid-generate} 之后，完整 SID 的 collision 已经处于很低水平，当前的主要瓶颈并不是未修复的全局碰撞，而是发生在语义前缀内部的局部歧义：prefix 往往已经正确，但最后一层 leaf token 仍然难以稳定预测。基于这一观察，我们提出 \textbf{MGR-SID}，一种建立在 MiniOneRec 语义 \rqvae{} backbone 之上的层级感知图正则化 tokenizer。当前的 \textbf{v1} 实例包含一个仅用训练数据构造的 graph bank，其中包括净化后的 coarse collaborative graph、spectral middle-resolution graph 以及净化后的 local transition graph，并在不同 SID level 上施加不同的图结构约束。首先，我们在 Industrial 和 Office 上通过 probe 证明：spectral middle-resolution graph 在 same-$l_2$ 歧义桶上的判别力显著强于第一版 heuristic middle view。随后，我们将 tokenizer 训练制度对齐到上游 MiniOneRec，并在公平设定下验证训练时 graph regularization。Industrial 上，层级图正则将最终生成 index 的 collision 从 $0.00353$ 降到 $0.00326$，并把测试目标落在多叶 same-$l_2$ bucket 的比例从 $68.94\%$ 降到 $61.31\%$。把这套 SID 推进到下游 SFT 后，我们得到一个更细致的结果：hierarchy-aware tokenizer 在短程排序上带来提升（$\mathrm{NDCG}@3: 0.0777 \rightarrow 0.0802$），同时减少了 evaluate 阶段的 same-prefix 错误，但在更广义的 $@10+$ 指标上还没有全面超过 semantic baseline。整体来看，这些结果说明，相比统一的 graph feature fusion，graph-structured collaboration 更适合作为 level-specific 的结构监督进入 SID 学习，但 tokenizer 到 SFT 的传递机制仍然只解决了一部分问题。
